\section{Introduction}\label{sec:introduction}

Variation of stability changes the birational geometry of moduli spaces
of sheaves on a K3 surface. Near a polystable sheaf, an Ext-quiver can
describe the analytic contraction and its exceptional strata. On the
locally free locus, the Hermitian--Einstein correspondence gives another
geometric structure: a tangent vector has the $L^2$ norm of its harmonic
Kodaira--Spencer representative. An analytic quiver description does not
identify this metric. In particular, it does not determine its behavior
when several exceptional strata approach one another at different rates.

We establish such a metric comparison for a rank-three wall crossing whose
local complex geometry is the stratified Grassmannian Mukai flop. The wall
contains an eight-dimensional stratum with projective-plane fibres and a
more reducible point with a Grassmannian fibre. The comparison holds on a
full neighborhood containing both strata. It gives relative asymptotics
for intrinsic distances, the complete metric obtained by rescaling the
contracting region, and a quantitative comparison with the surrounding
wall metric.

\subsection{The moduli problem and metric convention}

Section~\ref{sec:k3} constructs a smooth genus-four K3 surface $X$ containing
smooth rational curves $D_1,D_2$ such that
\begin{equation}\label{eq:intersection-intro}
D_1^2=D_2^2=-2,\qquad D_1D_2=5,\qquad H=D_1+D_2\ \text{is ample}.
\end{equation}
Put $L_i=\cO_X(D_i)$ and
\[
E_0=(L_1\otimes\CC^2)\oplus L_2,\qquad
v=(3,2D_1+D_2,0),\qquad D=D_1-D_2.
\]
We fix the holomorphic determinant $L_1^2L_2$. The Mukai vector $v$ is
primitive, $v^2=10$, and $c_2(E_0)=8$. At the polarization $H$, the
bundle $E_0$ is both slope-polystable and Gieseker-polystable.

Choose a smooth path of K\"ahler forms $\omega_t$ in $H+tD$, for $|t|$
small, with $\omega_0$ Ricci-flat, whose existence in $H$ follows from
Yau's theorem \cite{Yau78}. Let
$\cM_t^{\rm bun}$ be the moduli space of slope-polystable locally free
bundles of vector $v$ and the specified determinant. We use its
Hermitian--Yang--Mills connection topology. On each smooth orbit-type
stratum, where the conjugacy class of the effective compact stabilizer
is fixed, its metric $g_t$ is the harmonic Dolbeault metric, with
\begin{equation}\label{eq:dolbeault-norm}
\|a\|_{\omega_t,h}^2
=\int_X\tr(i a^\dagger\wedge a)\wedge\omega_t .
\end{equation}
There is no division by volume. The real connection variation $a-a^\dagger$
has squared norm twice~\eqref{eq:dolbeault-norm}. The determinant line has
its $\omega_t$-Hermitian--Einstein metric; changing a scalar bundle metric
does not affect the endomorphism norms in~\eqref{eq:dolbeault-norm}.

The length distance $d_t$ is defined using continuous paths that are
piecewise smooth on the smooth orbit-type strata, with local gauge
lifts. Finite-length limiting endpoints at lower strata are allowed.
Paths used to define $d_t$ may leave any prescribed coordinate
neighborhood; they remain in the locally free moduli component. We do
not impose this metric on non-locally-free sheaves.

\subsection{The comparison theorem}

We first specify the finite-dimensional quotient metric appearing in the result. Let
\[
\cV=\Hom(\CC^2,\CC^5)\oplus\Hom(\CC^5,\CC^2),\qquad x=(A,B),
\]
with Euclidean Hermitian norm $\|x\|^2=\tr A^*A+\tr BB^*$.
The group $G=GL_2(\CC)$ acts by
$(A,B)\mapsto(Ag^{-1},gB)$, with compact form $K=U(2)$. Define
\begin{equation}\label{eq:model-intro}
M_\tau=\{BA=0,\ A^*A-BB^*=\tau I_2\}/K,\qquad
\tau=\tau(t)=-\frac{28\pi}{3}t .
\end{equation}
Write $g_\tau^{\rm q}$ and $d_\tau^{\rm q}$ for its flat quotient metric
and length distance, and $r_\tau([A,B])=(\|A\|^2+\|B\|^2)^{1/2}$.
At $\tau=0$ these are understood on the stratified quotient.

The invariant $N=AB$ takes values in
\begin{equation}\label{eq:nilpotent-cone}
\cN=\{N\in\End(\CC^5):N^2=0,\ \rank N\le2\}.
\end{equation}
For $\tau>0$, $M_\tau=T^*\Gr(2,5)$; for $\tau<0$,
$M_\tau=T^*\Gr(3,5)$; and $M_0=\cN$. The quotient maps
$\pi_\tau:M_\tau\to\cN$ are the two incidence resolutions. Their fibres
are points over rank two, $\PP^2$ over rank one, and the full
Grassmannians over zero. Section~\ref{sec:model} recalls these
identifications with their metric normalizations.

\begin{theorem}\label{thm:main}
There are $R,t_0,C>0$ such that, for $|t|<t_0$, there are
homeomorphisms
\[
f_t:\{r_{\tau(t)}<R\}\subset M_{\tau(t)}
\longrightarrow U_t\subset\cM_t^{\rm bun}
\]
onto open subsets. They are smooth and biholomorphic on the corresponding
smooth complex orbit-type strata. The vertex maps to $[E_0]$ at
$t=0$, and the whole zero section is included for all sufficiently small
nonzero $t$. Every point of $U_t$ is slope-stable when $t\ne0$.

On every smooth piece and for every tangent vector,
\begin{equation}\label{eq:main-tensor}
(1-C(|t|+r_\tau^2))g_\tau^{\rm q}
\ \le\ f_t^*g_t\ \le\
(1+C(|t|+r_\tau^2))g_\tau^{\rm q}.
\end{equation}
The constants are independent of the positive singular values of $AB$
and of their ratios to each other and to $|t|$.

For some $c,C'>0$, if $r_\tau(p),r_\tau(q)\le\rho\le cR$, then
\begin{equation}\label{eq:main-distance}
(1-C'(|t|+\rho^2))d_\tau^{\rm q}(p,q)
\le d_t(f_t(p),f_t(q))
\le(1+C'(|t|+\rho^2))d_\tau^{\rm q}(p,q).
\end{equation}
The middle distance allows paths throughout the locally free moduli
component.
\end{theorem}

The parameter in~\eqref{eq:model-intro} follows from the Chern degrees and
the harmonic pairing; it is not a choice of scaling exponent. Since
$r_\tau^2\ge2|\tau|$, the error in~\eqref{eq:main-tensor} is also $O(r_\tau^2)$.
At the wall, the rank-one piece parametrizes bundles $L_1\oplus F$, where
$F$ is stable of Mukai vector $(2,H,0)$. It meets the more reducible
bundle $E_0$. Thus the uniformity includes a genuine change from effective
circle symmetry to effective $U(2)$ symmetry.

For the integral ample classes $mH\pm D$, with $m$ sufficiently large,
$c_1(E_0)(mH\pm D)=9m\mp7$ is coprime to three. Slope and Gieseker
stability therefore coincide for these invariants. The two chamber
neighborhoods occur in the locally free loci of the ordinary projective
sheaf moduli spaces. The contraction diagram asserted here is local
and analytic. It does not require a global morphism to the ordinary
$H$-Gieseker wall moduli space.

\subsection{Distances and rescaled geometry}

Let $s_1\ge s_2\ge0$ be the singular values of $N=AB$. For $\kappa>0$ set
\begin{equation}\label{eq:profile-intro}
b_i^2=\frac{\sqrt{\kappa^2+4s_i^2}-\kappa}{2},\qquad
F_\kappa(b)=\int_0^b
\sqrt{\frac{\kappa+2u^2}{\kappa+u^2}}\,du.
\end{equation}
For $\kappa=0$ put $b_i=\sqrt{s_i}$ and $F_0(b)=\sqrt2\,b$.
Take a slightly larger chart than the one in Theorem~\ref{thm:main},
so that closed radius subsets are defined. Let
\[
\cE_t=f_t(\pi_\tau^{-1}(0)),\qquad
\cD_t=f_t(\pi_\tau^{-1}\{\rank N\le1\}\cap\{r_\tau\le R\}).
\]
These are the deepest fibre and the local contracted locus; at $t=0$
they become the vertex and the closure of the reducible stratum.

\begin{theorem}\label{thm:distances}
For $x$ of radius $r=r_\tau(x)\le cR$, with $\kappa=|\tau|$,
\begin{align}
d_t(f_t(x),\cE_t)
&=(1+O(|t|+r^2))
\sqrt{F_\kappa(b_1)^2+F_\kappa(b_2)^2},\label{eq:distance-zero}\\
d_t(f_t(x),\cD_t)
&=(1+O(|t|+r^2))F_\kappa(b_2).\label{eq:distance-rank}
\end{align}
The errors are multiplicative and uniform in all the indicated scales.
At zero distance the formulas mean equality to zero.

The metric induced on $\cE_t$ is $|\tau|$ times the standard
Grassmannian quotient metric, with relative error $O(|t|)$.
On the $\PP^2$ fibre over a rank-one matrix of singular value $s>0$,
the induced metric is $|\tau|$ times the Hopf-quotient Fubini--Study
metric, with relative error $O(|t|+s)$.
\end{theorem}

In particular,
\[
d_0(f_0(N),\cD_0)=(1+O(r_0^2))\sqrt{2s_2}.
\]
For nonzero $\tau$, the exact model profile crosses from
$s_2/\sqrt{|\tau|}$ when $s_2\ll|\tau|$ to $\sqrt{2s_2}$ when
$|\tau|\ll s_2$. Formula~\eqref{eq:distance-rank} controls the
intermediate ratios and remains relative even if this distance is
arbitrarily small compared with $r$. Fibre metrics in the theorem are
induced metrics; their intrinsic distances are not identified with
unrestricted ambient distances between fibre points.

\begin{corollary}\label{cor:limits}
Fix a zero-section point in $M_{+1}$ or $M_{-1}$. On the corresponding
side of the wall, dividing $g_t$ by $\kappa=|\tau(t)|$ gives the complete
pointed length-space limit
$(M_{+1},g_{+1}^{\rm q})$ or $(M_{-1},g_{-1}^{\rm q})$.
More precisely, the maps
$\iota_t(\widehat x)=f_t(\sqrt\kappa\,\widehat x)$ satisfy
\begin{equation}\label{eq:inner-intro}
(1-C\kappa(1+\widehat r^2))g_{\pm1}^{\rm q}
\le\kappa^{-1}\iota_t^*g_t
\le(1+C\kappa(1+\widehat r^2))g_{\pm1}^{\rm q}.
\end{equation}
This holds on expanding domains $\widehat r\le R_t$ whenever
$R_t\to\infty$ and $\kappa R_t^2\to0$, and implies pointed
Gromov--Hausdorff convergence on every bounded pointed ball.

On the common rank-two locus, using $N$ to identify the chamber with
the wall, the actual chamber and wall metrics have relative difference
\begin{equation}\label{eq:outer-intro}
O\left(|t|+r_0^2+\frac{|\tau|}{s_2}\right)
\end{equation}
when $|\tau|/s_2$ is sufficiently small.
\end{corollary}

The abstract Riemannian spaces $M_{+1}$ and $M_{-1}$ are isometric
by adjoint interchange. Their contractions over the fixed matrix $N$
are different. The complete metrics in the corollary are classical;
the assertion identifies them as limits of the actual harmonic metric.
Neither the corollary nor the tensor estimate asserts smooth convergence
or curvature bounds.

\subsection{Relation to previous work and the analytic argument}

The complex geometry uses established K3 deformation theory. Budur--Zhang
\cite[Theorem~1.2]{BZ19} prove formality for a sheaf polystable with
respect to any ample polarization. Arbarello--Sacc\`a
\cite[Propositions~4.4 and~6.1]{AS18} give the equivariant analytic
quadratic germ and its quiver interpretation. Their adjacent-polarization
wall-crossing theorem concerns pure one-dimensional sheaves; it is not
used for the rank-three bundles here. Buchdahl--Schumacher
\cite{BS21,BS22} construct harmonic analytic deformation quotients.
Delloque's simultaneous deformation theory
\cite[Propositions~4.3--4.6 and Theorem~7.4]{Del25}
supplies corrected HYM equations and an open actual moduli neighborhood
on the integrable locus, without an abelian-stabilizer assumption.

Local metric comparison in reduction is also not new. Sjamaar
\cite[arXiv version, Lemma~2.13 and its proof]{Sja95} compares the reduced
K\"ahler metric with its flat holomorphic slice metric on the top stratum
near a chosen quotient point. Sjamaar--Lerman's singular-reduction theory
\cite{SL91}, using the classical symplectic normal form, and Mayrand's
complex Hamiltonian normal form \cite[Theorems~1.3 and~1.4]{May22}
preserve structures different from the full Riemannian metric. A real symplectic
coordinate change need not preserve an independently imposed complex
integrability equation. Namikawa--Odaka \cite[Remark~6.14]{NO26}
discuss the stronger metric comparison problem in the setting of
symplectic singularities; their isolated-singularity and finite
hyperk\"ahler quotient results do not directly give the simultaneous
estimate in Theorem~\ref{thm:main}.

The complete cotangent-Grassmannian metric and its quotient description
are classical \cite{BG96}. Qualitative Gromov--Hausdorff continuity
through suitable projective Calabi--Yau flops is proved by Rong--Zhang
\cite[Theorem~1.2]{RZ11}. Actual Weil--Petersson forms on bundle moduli
and their extensions are studied in \cite{BS21,BS23}; explicit metrics
and degenerations for framed bundles on the blown-up plane appear in
\cite{BS26}. Gauge slices and reduction by the residual finite-dimensional
group are part of the earlier Yang--Mills moduli theory
\cite{GP89}. We do not claim these constructions or the finite
Grassmannian scaling as new.

The additional issue is to compare the \emph{actual} real moment levels
on the exact obstructed complex locus, with an error independent of the
smallest positive infinitesimal gauge direction. A usual implicit-function
bound can lose that direction. We replace it by a projection estimate.
If $\mu=A^*A-BB^*$ and $\rho$ is the centered corrected moment map,
then for every orthogonal projection $P$ on $\CC^2$,
\[
|\tr(P(\rho-\mu))|
\le\epsilon\bigl(\|AP\|^2+\|PB\|^2\bigr).
\]
The right side is the squared norm of the circle action generated by
$P$. Testing this estimate on spectral projections of the positive
Hermitian factor of a complex gauge transformation bounds that factor
uniformly through rank loss. Projection-relative curvature estimates and
finite-dimensional balancing by spectral projections were developed for
repeated stable factors in \cite{Wue26a,Wue26c}. Here we adapt this method
to the exact obstructed K3 deformation locus and compare different real
moment levels. Section~\ref{sec:balancing} gives the required argument
in full.

The harmonic norm requires a separate identification. We first minimize
over the uniformly coercive complement of the parallel gauge algebra,
as in the complementary-gauge method of \cite{Wue26b}. The result is a
smooth positive Hermitian form on the full deformation space. A central
automorphism acting by minus the identity removes its linear term.
Only then do we minimize over the remaining complex orbit directions.
Derivatives of the real balancing transformation are exact gauge terms
and disappear from that infimum. Finally, a radius estimate excludes
shorter paths leaving the chart. These steps give the actual intrinsic
conclusions, rather than merely a comparison of norms on an assumed
common locus.

The result is local and concerns locally free bundles. It gives no
global completion theorem, no bubbling description, and no general
metric normal-form theorem for hyperk\"ahler singularities.

The polynomial calculations of Section~\ref{sec:k3} and the distance
profile in Figure~\ref{fig:distance} can be reproduced from the
source accompanying this paper at
\url{https://github.com/bwuebben/stability-degeneration}.
